\section{Three connected research questions}
In a limited concert-ticket market, a consumer must choose between buying immediately to secure a ticket and waiting for a possible discount while risking sell-out. This creates a strategic scarcity problem because the value of waiting depends not only on the consumer's own choice but also on competition from other buyers for limited supply. My Intellectual Statement II introduced AI advice into this setting: \emph{greater reliance need not produce better decisions}. Consumers, ticket platforms, and consumer-protection educators therefore need a benchmark for distinguishing justified reliance from advice-following that only appears successful after the outcome is known. Simon's bounded rationality~\cite{simon1955} motivates comparing feasible human reasoning with an ideal decision benchmark. Prior studies examine costly reliance~\cite{klingbeil2024}, trust dynamics~\cite{li2023}, and individual differences in AI-assisted judgment~\cite{pearson2026}, but they do not jointly connect market scarcity, advice reliability, and the welfare consequences of reliance in this ticket setting.

\textbf{Q1---Economics:} When does following AI advice increase a consumer's expected surplus, and how does this benchmark change as competition for limited tickets becomes more severe? \textbf{Q2---Computation:} Can an inspectable simulator distinguish simple advice-following from posterior-optimal choice and quantify the resulting regret? \textbf{Q3---Behavior:} Does observed reliance primarily reflect trust in the AI adviser or consumers' beliefs about ticket sell-out risk? The integrated question is when AI reliance is economically justified once market competition, advice reliability, and human beliefs are considered together. Figure~\ref{fig:teaser} illustrates how these three components connect.

\section{Answer to the economics question}

The economic setting is a scarce-ticket allocation problem in which multiple consumers compete for a limited supply $Q$. Each consumer can buy immediately ($B$) or wait ($W$) for a 10\% discount, but waiting creates a risk that tickets sell out before purchase. The full strategic problem would make this risk depend endogenously on other consumers' choices. The current PS1 benchmark instead holds aggregate competition fixed at $N$ and solves the focal consumer's Bayesian best response. Thus, the computation is a tractable decision benchmark within a broader strategic environment, rather than a market-equilibrium estimate.

The focal consumer observes price $P$, value $V$, supply $Q$, competition $N$, and advice $A$ of stated reliability $r$. Other buyers enter through the exogenous sell-out probability
$p=\min\{0.98,N/(N+8Q+1)\}$ inherited from my game. The consumer observes the advice, chooses $B$ or $W$, and then learns whether tickets sold out ($S=1$). A strategy maps the observed market inputs and advice into an action; $B$ or $W$ is one realized action. Assume $V>P>0$ and risk neutrality. Buying guarantees payoff $V-P$, while waiting gives zero if sold out and $V-0.9P$ otherwise.

The current solution benchmark is Bayesian expected-utility maximization~\cite{savage1972}. The stochastic extension assumes
$\Pr(A=B\mid S=1)=\Pr(A=W\mid S=0)=r$,
known to the consumer. For $A=W$,
\[
p_W=
\frac{p(1-r)}
{p(1-r)+(1-p)r},
\qquad
B \text{ is optimal iff }
p_A \geq
\frac{0.1P}{V-0.9P}.
\]
Here $p_A$ is the posterior sell-out risk; without advice the benchmark uses $p$. Therefore, even relatively reliable ``Wait'' advice can rationally be rejected when scarcity is sufficiently severe. This is a conditional mathematical prediction under the model, not an estimated market effect. The sell-out function and symmetric signal structure are assumed rather than measured. Endogenous competition among buyers, seller incentives, and heterogeneous risk preferences remain unresolved and motivate the next strategic extension.

\section{Answer to the computational question}

The notebook implements the computational component of Figure~\ref{fig:teaser} by exactly enumerating the two market states and two advice signals for each condition. Inputs are price $P$, value $V$, supply $Q$, competition $N$, and advice reliability $r$. The baseline policy always follows the advice; the main modification chooses the action with the larger posterior expected payoff, while ignoring advice provides an additional comparator. The primary metrics are expected surplus and regret relative to the posterior-optimal policy.

The smallest experiment fixes $P=150$, $V=220$, $Q=40$ and compares the three policies; a 25-condition grid varies $N$ and $r$. At $N=500,r=.8$, always-follow yields expected surplus $66.1657$ versus $70$ for Bayes (regret $3.8343$); a seeded simulation gives $66.202$. These are synthetic, not human, results; tests appear in Appendix~A. The game shows retrospective outcomes, whereas the notebook evaluates ex-ante payoffs.

These are obtained synthetic results, not evidence about human behavior. Threshold, signal-boundary, and policy-dominance checks pass; detailed outputs and tests are reported in Appendix~A. The existing interactive game uses deterministic scenario hashes and retrospective correctness, whereas the notebook replaces them with a stochastic signal model and ex-ante expected-payoff evaluation. Thus, the game illustrates the decision problem, while the notebook provides the reproducible benchmark used to evaluate it.

\section{Answer to the behavioral question}

My provisional behavioral mechanism is trust-driven acceptance of AI advice; the competing explanation is that consumers who follow the same advice hold higher subjective beliefs about ticket sell-out risk. Prior work on trust dynamics~\cite{li2023} motivates modeling reliance across repeated interactions, while Pearson et al.~\cite{pearson2026} show that attitudes toward AI and decision performance need not coincide. The existing game outputs cannot distinguish these mechanisms.

A feasible pilot would hold $P$, $V$, and $Q$ fixed, randomize low/high competition $N$ and advice reliability $r\in\{.6,.8\}$, elicit perceived sell-out probability before and after advice, and record choices and expected regret. A second comparison would vary displayed competition while holding the supplied sell-out probability constant. If reliance changes mainly with stated AI reliability after accounting for elicited risk beliefs, the evidence would support a trust-based mechanism. If reliance is instead explained primarily by changes in sell-out beliefs, I would weaken the trust account and revise the behavioral model accordingly.

The test cannot fully identify trust because belief reports may be noisy and residual choice differences may reflect misunderstanding or risk preferences; comprehension checks and a risk-preference measure are therefore needed. These behavioral mechanisms also affect the economic and computational design: subjective beliefs can replace the benchmark probability $p$ in the decision rule, allowing observed choices to be compared with the normative benchmark. Following my September 7 observation of Session C during the workshop, a later repeated-round experiment would randomize advice-success histories and test how reliance and beliefs update after experience. Human recruitment, sample size, and causal estimates remain planned.

\section{Advanced development and future directions}

The enduring question is when people should delegate judgment under uncertainty. Simon's bounded rationality and 1978 Nobel recognition~\cite{simon1955,nobel1978}, together with Newell and Simon's 1975 Turing Award~\cite{newell1976,turing1975}, provide the foundation. Today's AI recommendations embed delegation directly in scarce decisions. The unresolved boundary is whether reliance reflects environmental-risk or AI-quality beliefs, and the welfare loss from mistaken reliance.

Economically, I define an expected-surplus benchmark; computationally, I make it inspectable and reproducible; behaviorally, I propose separating trust from market-risk beliefs. The next study tests these mechanisms with humans; longer-term work endogenizes strategic competition and repeated learning.

\textbf{Expected 2056 Nobel Prize and Turing Award contribution statement.}
By 2056, I aspire to develop an integrated economic, computational, and behavioral framework identifying when AI delegation improves decisions under uncertainty. Building on Simon and Newell, the framework would combine welfare benchmarks, reproducible computation, and validated behavioral mechanisms to distinguish beliefs about the environment from beliefs about AI quality across consequential settings.

\subsection*{Open Science Statement}

The project materials are available in the \GitHubLink, with an executable version in the \ColabLink. The repository contains the ticket-advice model, executable notebook, tests, synthetic outputs, and preserved original game. The model uses only the Python standard library; from the repository root, run \texttt{python3 research/model.py}, or open the notebook in Colab and select ``Run all.'' The submitted version is identified by commit \texttt{\CommitHash}. Reuse terms and detailed reproducibility records are reported in Appendix~A.

\subsection*{Statement of Contribution to the UN's SDGs}
The interactive learning tool is available through the \DemoLink. The open-access game supports a proposed \textbf{SDG 4---Quality Education}~\cite{sdg4} lesson for introductory students: predict a choice, change competition or reliability, then compare outcome feedback with the notebook's expected-payoff benchmark. Access requires a browser and internet; the static game needs no API key. A pre/post explanation task could assess whether students distinguish good decisions from lucky outcomes. The resource exists; learning gains remain untested.
